% !TEX root = ../main.tex

\section{研究方法}

\subsection{时空深度学习模型}

\subsubsection{U-TAE模型}

% 待重写

\subsubsection{3D U-Net模型}

% 待重写

\subsubsection{ConvGRU模型}

% 待重写

\subsection{实验设计}

本研究围绕玉米、大豆与其它三类目标，设计站内分类实验与跨站点空间泛化实验，以评估三种时空深度学习模型的分类能力与空间泛化性能。

\subsubsection{站内分类实验}

在站内实验中，每个站点使用2023年多时相HLS影像和CDL标签构建的400个非重叠图块样本，采用5折交叉验证策略进行模型训练与评估。具体而言，400个样本被划分为5个互不重叠的子集，每轮实验选取其中3个子集用于训练、1个子集用于验证（早停判断）、1个子集用于测试。重复5轮后，每个样本恰好被测试一次，形成完整的站内Out-of-Fold（OOF）预测。

站内结果从两个层面进行汇报。其一，将5轮测试各自计算的分类指标取均值与标准差，反映模型在不同训练-测试划分下的性能波动；其二，将5轮测试的混淆矩阵逐元素汇总后重新计算各项指标，得到基于完整OOF预测的总体精度，用于后续与跨站点结果进行一致口径的比较。

三种模型均基于PyTorch框架实现，采用统一的训练策略与超参数设置：优化器为AdamW，损失函数为交叉熵损失，初始学习率设为0.0001，批次大小为4，使用余弦退火学习率调度策略（$T_{\max}=50$，$\eta_{\min}=0.00005$），权重衰减为0.00005，梯度裁剪阈值为5.0，并引入早停机制（容忍轮次为15，最小改善量为0.0001）以防止过拟合。训练阶段对每个原始图块施加$0^{\circ}$、$90^{\circ}$、$180^{\circ}$和$270^{\circ}$四种旋转增强，使每折的有效训练样本量从240扩展至960。所有模型均在相同硬件环境下进行训练。

\subsubsection{跨站点空间泛化实验}

在跨站点泛化实验中，依次选取四个站点之一作为源站点，利用其站内实验训练所得的5个fold模型，分别对其余三个目标站点的全部400个图块进行逐图块推理。每条迁移路径产生5次独立的完整预测，对每次预测将400个图块的混淆矩阵汇总后计算各项指标，再对5次结果取均值与标准差。迁移过程中不对模型进行微调或集成，以评估模型结构本身的空间泛化能力。4个站点轮换作为源站点，共构成12条有向迁移路径，对应3种模型共36组结果、180次推理。

为量化跨站点应用的精度损失，本文以同一模型在目标站点的站内OOF汇总精度为参照基准，定义泛化退化指标（以玉米F1为例）：
\begin{equation}
\Delta\text{F1}_{\text{corn}} = \text{F1}_{\text{corn,within}} - \text{F1}_{\text{corn,cross}}
\end{equation}
其中$\text{F1}_{\text{corn,within}}$为该模型在目标站点的站内OOF汇总玉米F1分数，$\text{F1}_{\text{corn,cross}}$为源站点模型在目标站点的跨站点玉米F1五折均值，正值表示精度下降。大豆泛化退化$\Delta\text{F1}_{\text{soybean}}$同理计算。

\subsection{精度评价指标}

模型性能评估基于预测结果与像素级真实标签构建的混淆矩阵。对于三分类问题（玉米、大豆、其他），混淆矩阵元素$C_{ij}$表示真实类别为$i$但被预测为类别$j$的像元数。基于混淆矩阵，各类别F1分数计算如下：
\begin{equation}
\text{Precision}_i = \frac{C_{ii}}{\sum_{j} C_{ji}}, \quad
\text{Recall}_i = \frac{C_{ii}}{\sum_{j} C_{ij}}, \quad
\text{F1}_i = \frac{2 \cdot \text{Precision}_i \cdot \text{Recall}_i}{\text{Precision}_i + \text{Recall}_i}
\end{equation}
其中$i$为类别索引。总体精度（OA）定义为：
\begin{equation}
\text{OA} = \frac{\sum_{i} C_{ii}}{\sum_{i,j} C_{ij}}
\end{equation}

由于本研究的核心关注点是玉米和大豆的识别与泛化能力，本文分别报告玉米F1（$\text{F1}_{\text{corn}}$）和大豆F1（$\text{F1}_{\text{soybean}}$）作为主要评价指标。同时计算OA作为总体性能参考。

为检验不同模型的跨区域泛化性能是否存在统计学上的显著差异，本文将12条迁移路径作为配对单位，分别对三种模型的玉米F1和大豆F1进行Friedman非参数检验，并在总体差异显著时采用Wilcoxon符号秩检验进行两两比较，使用Holm方法校正多重检验。

\subsection{硬件和软件配置}

建模过程是在Linux工作站（Ubuntu 22.04）上进行的，该工作站具有一个Intel Xeon Gold 6430处理器（2.1~GHz/32 Core/60~M），120~GB RAM和1个NVIDIA GeForce RTX 4090显卡（24~GB RAM）。所有模型都使用PyTorch库实现。
